\documentclass{article}

\usepackage{neurips_2026}

\usepackage[utf8]{inputenc}
\usepackage[T2A,T1]{fontenc}
\usepackage[english,russian]{babel}
\usepackage{tempora} % Times-compatible font with T2A Cyrillic
\usepackage{hyperref}
\usepackage{url}
\usepackage{booktabs}
\usepackage{amsfonts}
\usepackage{amssymb}
\usepackage{amsmath}
\usepackage{nicefrac}
\usepackage{microtype}
\usepackage{xcolor}
\usepackage{graphicx}
\usepackage{array}
\usepackage{enumitem}
\usepackage{placeins}

\newcommand{\method}[1]{\texttt{#1}}

\title{Коллективная семантическая память: слияние, доверие и устаревание для пиринговой памяти в многоагентной навигации}

%\author{%
%  Anonymous Authors \\
%  Anonymous Institution \\
%  Anonymous Address \\
%  \texttt{anonymous@anonymous.edu} \\
%}

\begin{document}

\maketitle

\begin{abstract}
\textbf{Проблема.} Коллективную память часто считают монотонной: чем больше общей информации, тем лучше. В кооперативной навигации с дефицитными семантическими целями эта интуиция не работает. Быстрое глобальное распространение информации улучшает материализацию целей, но заставляет агентов состязаться за одни и те же цели; медленный обмен избегает конкуренции, но оставляет агентов с устаревшей или неполной памятью.
\textbf{Подход.} Мы реализуем четыре топологии памяти над одним и тем же символьно-семантическим субстратом: независимые приватные памяти, общую шину событий, центральный консенсусный агрегатор и пиринговую широковещательную рассылку с каденцией (cadence) $K$. Затем мы вводим минимальную \emph{коллективную семантическую память} (CSM): пиринговую рассылку плюс явные правила \emph{слияния} (merge), \emph{доверия} (trust) и \emph{устаревания} (staleness). Показатели стадий Q/R/M/C используются только как измерительная инструментовка, а не как вклад данной статьи.
\textbf{Результаты.} На многоагентной развёртке из 35\,640 запусков семейство пиринговых протоколов с фиксированной каденцией демонстрирует сдвиг узкого места: быстрая рассылка насыщает материализацию, но обрушивает завершение, тогда как медленная рассылка имеет противоположный профиль. Среднее время до успеха имеет внутренний оптимум при $K{=}8$. На бенчмарке дефицита из 3\,240 запусков минимальная CSM строго доминирует над пятью пиринговыми базовыми вариантами с фиксированной каденцией по доле успеха: CSM достигает $0.617$ при 95\%-м ДИ $[0.610,0.624]$, выше верхней границы ДИ каждого peer-$K$.
\textbf{Границы применимости.} Эксперименты проведены на грид-мирах. Утверждение состоит не в том, что это полная теория коллективной памяти; это воспроизводимое архитектурное разбиение, показывающее, что слияние/доверие/устаревание --- правильные явные поверхности управления после того, как одна лишь каденция исчерпана.
\end{abstract}

\section{Введение}

Многие многоагентные системы обмениваются информацией либо записывая
в единую глобальную память, либо посылая обучаемые сообщения, либо
периодически обмениваясь локальным состоянием. Эти механизмы обычно
оцениваются сквозным образом. Если успех растёт, коммуникация
признаётся полезной; если нет --- причина остаётся непрозрачной. В этой
статье изучается более узкий архитектурный вопрос: когда несколько
агентов поддерживают семантическую память о местах, чем именно следует
делиться, с какой частотой и с каким весом?

Ответ не сводится к <<делиться всем и как можно скорее>>. В условиях
дефицита, когда $N$ агентов конкурируют за $T<N$ семантических целей,
свежая глобальная память может привести к тому, что все агенты выберут
одну и ту же цель. Память при этом точна, но бесполезна. И наоборот,
отсутствие коммуникации предотвращает стадное поведение, но лишает
агентов полезных открытий, сделанных соседями. Задача проектирования
--- это, следовательно, компромисс между семантической свежестью и
конкуренцией за занятость целей.

Мы делаем три вклада.
\begin{enumerate}[leftmargin=20pt,itemsep=2pt,topsep=2pt]
\item \textbf{Архитектурное разбиение.} Мы реализуем и сравниваем четыре топологии памяти над одним и тем же символьным субстратом: независимую, общую шину, центральный агрегатор и пиринговую широковещательную рассылку.
\item \textbf{Результат о каденции.} Мы показываем, что пиринговая рассылка обнажает реальный координационный параметр $K$: быстрая рассылка улучшает $P(M^\star|R^\star)$, но вредит $P(C^\star|M^\star)$; метрика эффективности имеет внутренний оптимум.
\item \textbf{Минимальная CSM.} Мы реализуем коллективную семантическую память с тремя явными правилами: слияние, доверие и устаревание. Минимальная CSM достигает новой точки на плоскости M$\times$C и строго доминирует над пиринговыми базовыми вариантами с фиксированной каденцией по доле успеха на бенчмарке дефицита.
\end{enumerate}

\paragraph{Связь с сопутствующей статьёй о Q/R/M/C.}
Сопутствующая статья рассматривает Q/R/M/C как диагностический
фреймворк. Здесь эти показатели стадий используются как измерительные
инструменты. Основной объект иной: само устройство распределённой
символьной памяти.

\section{Субстрат памяти и задача}
\label{sec:substrate}

Каждый агент наблюдает грид-среду и порождает семантические
свидетельства о местах, такие как \texttt{water\_source},
\texttt{hazard\_edge} и \texttt{open}. Локальная память хранит записи
о местах, уверенности в тегах, счётчики посещений, свежесть и
статистику переходов. Запросы планировщика извлекают
места-кандидаты путём сопоставления обязательных, предпочтительных и
штрафных тегов. Выбранный кандидат материализуется в конкретную клетку
и исполняется локальным контроллером.

Многоагентная задача использует сетку $12{\times}10$. Каждый агент
должен достичь клетки с тегом \texttt{water\_source}. Цели дефицитны:
$N{\in}\{3,5,8\}$ агентов, $T{\le}N$ целей, три планировки, три
плотности опасностей и 40 сидов для основной развёртки по каденции.
Каждый запуск логирует события Q/R/M/C по каждому агенту:
\[
Q^\star \rightarrow R^\star \rightarrow M^\star \rightarrow C^\star,
\]
где $M^\star$ означает, что семантический кандидат был зафиксирован
как исполнимая цель, а $C^\star$ --- что цель была достигнута. Эти логи
не являются архитектурным вкладом; они позволяют увидеть, где именно
отказывает каждая топология.

\section{Четыре топологии коллективной памяти}
\label{sec:archs}

\begin{figure}[t]
  \centering
  \includegraphics[width=\linewidth]{figures/fig_architectures.pdf}
  \caption{Четыре топологии памяти. \textbf{(a)} Независимая: приватная память у каждого агента, обмена нет. \textbf{(b)} Общая шина: одна глобальная память, каждый агент читает и пишет на каждом такте. \textbf{(c)} Центральный агрегатор: приватные памяти плюс хаб, сливающий снимки в консенсусный отчёт. \textbf{(d)} Пиринговая рассылка: приватные памяти с обменом снимками каждые $K$ тактов и слитыми представлениями у каждого агента.}
  \label{fig:architectures}
\end{figure}

Четыре топологии используют одну и ту же локальную память и
планировщик. Они различаются только тем, где происходит слияние.

\begin{description}[leftmargin=10pt,topsep=2pt,itemsep=2pt]
\item[Независимая.] Каждый агент владеет приватным хранилищем событий и графом концептов. Канала коммуникации нет.
\item[Общая шина.] Все наблюдения дописываются в единую глобальную шину событий. Один общий граф концептов перестраивается из объединения всех событий.
\item[Центральный агрегатор.] Каждый агент владеет приватным графом концептов; центральный консенсусный слой забирает снимки и вычисляет один взвешенный по доверию отчёт для всех.
\item[Пиринговая рассылка.] Каждый агент владеет приватным графом концептов и приватным слитым представлением. Каждые $K$ тактов агенты рассылают снимки; каждый агент сливает свою память с последним снимком каждого соседа.
\end{description}

Ключевое различие --- не в примитиве кластеризации. Оно в том, на
каком уровне этот примитив инстанцируется:
\[
\begin{array}{rcl}
\text{независимая} &:& N \text{ графов, }0\text{ слияний}\\
\text{общая шина} &:& 1 \text{ граф, }1\text{ неявное глобальное слияние}\\
\text{центральная} &:& N \text{ графов }+1\text{ центральное слияние}\\
\text{пиринговая} &:& N \text{ графов }+N\text{ приватных слияний.}
\end{array}
\]

\section{Каденция обнажает механизм отказа}
\label{sec:cadence}

У пиринговой топологии есть один явный параметр: каденция рассылки $K$.
Основная развёртка перебирает четыре архитектуры и восемь пиринговых
каденций $K{\in}\{1,2,4,8,16,32,48,64\}$, всего 35\,640 запусков.
Рисунок~\ref{fig:conditional-rates} суммирует результат.

\begin{figure}[t]
  \centering
  \includegraphics[width=\linewidth]{figures/fig_conditional_rates.pdf}
  \caption{Условные показатели и среднее время до успеха по развёртке из 35\,640 запусков. Быстрый обмен улучшает материализацию целей, но вредит завершению из-за конкуренции; медленный обмен имеет противоположный профиль. Среднее время до успеха имеет внутренний оптимум при $K{=}8$.}
  \label{fig:conditional-rates}
\end{figure}

Быстрая пиринговая рассылка и централизованный обмен живут в режиме
насыщения M и коллапса C: агенты надёжно фиксируют цели, но многие
фиксации указывают на цели, которые преследуют и соседи. Независимая
топология и медленная пиринговая рассылка живут в режиме дефицита M и
насыщения C: меньше агентов фиксируются на хороших целях, но те, кто
фиксируется, реже сталкиваются с занятостью со стороны соседей.

\begin{figure}[t]
  \centering
  \includegraphics[width=0.85\linewidth]{figures/fig_bottleneck_shift.pdf}
  \caption{Сдвиг узкого места по пиринговым каденциям. $P(M^\star|R^\star)$ падает, а $P(C^\star|M^\star)$ растёт с ростом $K$. Практическая метрика эффективности, среднее время до успеха, минимизируется при $K{=}8$.}
  \label{fig:bottleneck-shift}
\end{figure}

По всем пиринговым запускам корреляция Спирмена $P(M^\star|R^\star)$ с
$K$ равна $-0.608$, а $P(C^\star|M^\star)$ с $K$ равна $+0.420$ (обе с
$p<10^{-4}$). Более сильная сигнатура видна при фиксированной каденции:
внутрикаденционная корреляция Пирсона между звеньями M и C достигает
пика $-0.61$ при $K{=}4$ и затухает к шуму с ростом $K$. Каденция,
следовательно, не просто настраивает <<объём коммуникации>>. Она
перемещает узкое место между семантическим принятием обязательства и
завершением на разделяемом ресурсе.

\section{Минимальная коллективная семантическая память}
\label{sec:csm}

Развёртка по каденции показывает, что обмен снимками с фиксированной
частотой обнажает лишь одну поверхность управления. Коллективная
семантическая память должна предоставлять как минимум три:
\begin{description}[leftmargin=14pt,itemsep=1pt,topsep=2pt]
\item[Слияние.] Как комбинируются свидетельства соседей об одном и том же семантическом месте.
\item[Доверие.] Как взвешиваются свидетельства разных соседей, когда они конфликтуют или становятся ненадёжными.
\item[Устаревание.] Как старые свидетельства затухают, прежде чем повлиять на извлечение.
\end{description}

Минимальная CSM использует пиринговую рассылку при $K{=}8$ и
накладывает поверх три правила:
\begin{description}[leftmargin=14pt,itemsep=1pt,topsep=2pt]
\item[\textit{Слияние.}] Оценка места равна собственному свидетельству с весом 1 плюс свидетельства соседей с весом $\mathrm{trust}_i \cdot e^{-\alpha \cdot \mathrm{age}}$, где $\alpha=0.05$. Порог включения $\tau=0.30$.
\item[\textit{Доверие.}] Экспоненциальное скользящее среднее (EMA) по каждому соседу, отражающее согласованность успеха извлечения. Соседи, чей последний снимок поддерживает локально зафиксированную цель, наращивают доверие; остальные затухают к $0.4$. Скорость обновления $\beta=0.10$.
\item[\textit{Устаревание.}] Каждый широковещательный снимок снабжён меткой времени; вес слияния затухает как $e^{-\alpha\cdot\mathrm{age}}$. Снимки старше $10K$ тактов вычищаются.
\end{description}

Реализация является полностью совместимой заменой пиринговой памяти с
фиксированной каденцией:
\method{experiments/collective\_semantic\_memory/csm\_memory.py}.

\subsection{Результаты}

Мы сравниваем CSM с пиринговой рассылкой с фиксированной каденцией при
$K \in \{1,4,8,16,64\}$ на ячейках дефицита
$(N,T){\in}\{(3,2),(5,3),(8,5)\}$, трёх планировках, трёх плотностях
опасностей и 20 сидах: всего 3\,240 запусков.

\begin{table}[ht]
\caption{Минимальная CSM против пиринговой рассылки с фиксированной каденцией на бенчмарке дефицита. Нижняя граница 95\%-го бутстреп-ДИ CSM превышает верхнюю границу ДИ каждого peer-$K$.}
\label{tab:csm-results}
\centering
\small
\begin{tabular}{l c c c}
\toprule
Архитектура & $P(M^\star|R^\star)$ & $P(C^\star|M^\star)$ & Доля успеха (95\% ДИ) \\
\midrule
peer ($K{=}1$)  & 0.991 & 0.586 & 0.577 [0.572, 0.586] \\
peer ($K{=}4$)  & 0.955 & 0.623 & 0.581 [0.572, 0.590] \\
peer ($K{=}8$)  & 0.730 & 0.871 & 0.599 [0.591, 0.607] \\
peer ($K{=}16$) & 0.697 & 0.892 & 0.597 [0.589, 0.606] \\
peer ($K{=}64$) & 0.684 & 0.900 & 0.601 [0.594, 0.609] \\
\textbf{CSM (минимальная)} & \textbf{0.782} & \textbf{0.825} & \textbf{0.617 [0.610, 0.624]} \\
\bottomrule
\end{tabular}
\end{table}

CSM занимает третий режим. Быстрые рассылки обладают превосходной
материализацией, но плохим завершением; медленные --- наоборот. CSM
удерживает оба условных звена выше 0.78. Абсолютный прирост успеха
скромен, от $+1.7$ до $+4.0$ процентных пунктов, но стадийная
декомпозиция показывает, почему он содержателен: архитектура меняет
совместную рабочую точку M$\times$C, а не просто перемещается вдоль
кривой фиксированной каденции.

\paragraph{Почему CSM выигрывает: вентиль, а не обмен.} Провенансная
аннотация тех же запусков даёт механизм. CSM обменивается не
\emph{лучше}, чем пиринговые варианты с фиксированной каденцией, ---
она \emph{отказывается действовать по устаревшим или недоверенным
чужим свидетельствам}. Правило включения из раздела~\ref{sec:csm}
влечёт замкнутую формулу горизонта, за которым свидетельство соседа
уже не может сформировать намерение,
$\mathrm{age}_{\max}(\mathrm{trust}_i) = \alpha^{-1}
\ln(\mathrm{trust}_i \cdot c/\tau)$, и этот горизонт эмпирически
точен: измеренные пороги формирования намерений совпадают с формулой
такт в такт по всем уровням доверия (23/23.05, 18/18.59, 12/12.84 при
$\mathrm{trust}=1.0/0.8/0.6$), включая контрольную проверку по
зарегистрированному предсказанию на шести никогда не запускавшихся
ячейках ($6/6$ точных целочисленных точек излома) и дискретное
переключение доверия: $\mathrm{trust}\le\tau/c$ полностью обнуляет
обязательства по чужим свидетельствам. У пиринговых вариантов с
фиксированной каденцией такого вентиля нет --- их чужие свидетельства
никогда не истекают, --- и именно поэтому их завершение обрушивается
при быстром обмене. Полный анализ, обусловленный провенансом,
приводится в сопутствующей рукописи; здесь он служит механистическим
прочтением таблицы~\ref{tab:csm-results}.

\section{Тождество мест из смысла: устранение общей системы координат}
\label{sec:semantic-identity}

Всё изложенное выше --- и категорное прочтение ниже --- предполагает
общий пространственный базис $X$: <<то же место>> означает <<те же
$(x,y)$ в системе координат, которую среда выдаёт каждому агенту>>.
Этот раздел устраняет это предположение и сообщает, что оно скрывало.
Каждый агент теперь держит \emph{приватную} систему координат
(истинную сетку под секретным сдвигом, а в последнем эксперименте ---
и под секретным поворотом, кратным $90^\circ$); широковещательные
снимки несут приватные координаты отправителя, бессмысленные для
получателя, пока тот не решит, что места отправителя \emph{значат}.

\paragraph{Механизм.} \emph{Отпечаток места} (place fingerprint) ---
это локальная семантическая констелляция вокруг посещённой клетки:
какие содержательные теги (опасности, стены, вода) находятся на каких
относительных смещениях --- инвариантно к сдвигу и без координат. Два
агента сопоставляют места по \emph{взаимно уникальному} соответствию
отпечатков (каждая сторона --- единственный кандидат другой;
однонаправленная уникальность вклеивает далёкие подобия в окрестность
получателя, а узоры, состоящие только из границ, исключаются как
описания формы мира, а не места). Межагентная система координат ---
это модальный сдвиг по $\ge 3$ согласующимся сопоставленным парам; при
поворотах по одной гипотезе на каждые $90^\circ$ переиспользуется тот
же выравниватель, и победитель обязан превосходить консенсус
ближайшего конкурента в $2\times$. Чужие свидетельства --- включая
места, которых получатель никогда не видел, --- затем переносятся
через восстановленную систему координат.

\paragraph{Результаты.} Зарегистрированы до выполнения; полный
протокол в сопутствующей рукописи.
\begin{enumerate}[leftmargin=20pt,itemsep=2pt,topsep=2pt]
\item \textbf{Точное восстановление.} Свидетель--путник: секретный
  сдвиг восстанавливается точно в $100\%$ запусков в мирах с
  ориентирами, а при полных $SE(2)$-системах координат --- в $80/80$
  комбинаций поворот$\times$сдвиг$\times$сид, со $100\%$-м завершением
  на фиксациях, принятых по чужим свидетельствам.
\item \textbf{Асимметрия отказов.} Там, где идентификация невозможна
  --- в безликом регионе или в $180^\circ$-симметричной планировке,
  где система координат невосстановима \emph{в принципе}, ---
  семантическое тождество не делает \emph{ни одного} обязательства
  (отказывает закрытым образом, fail closed), тогда как координатное
  тождество продолжает уверенно фиксировать неверные клетки ($100\%$
  неверных фиксаций; отказывает открытым образом, fail open).
\item \textbf{Полный стек.} При $N{=}4$ агентах в различных секретных
  системах координат выравнивание складывается онлайн (возникновение в
  $100\%$ эпизодов, средний такт ${\approx}15$) и восстанавливает
  $88\%$ разрыва производительности с общей системой координат (успех
  $0.444$ против $0.456$ у оракула против $0.350$ у координатного
  варианта); при случайных $SE(2)$-системах каждое семантическое
  обязательство указывает на настоящую воду ($8/8$), тогда как
  координатный стек производит $992$ фантомные фиксации ($0$
  настоящих).
\item \textbf{Вентиль ортогонален.} Повторный прогон анализа
  $\mathrm{age}_{\max}$ выше с правилом включения CSM ПОВЕРХ
  семантического тождества, при сдвиге системы координат и при
  повороте на $90^\circ$, воспроизводит ТЕ ЖЕ целочисленные точки
  излома ($20/14/9$; $\mathrm{bp}=20$ при повороте) --- вентиль
  доверие$\times$устаревание и восстановление системы координат суть
  независимые слои.
\end{enumerate}

Более сильное утверждение, обещанное в ограничениях ранних версий,
тем самым конструктивно исполнено: тождество мест (place identity)
может нестись одним лишь смыслом, оно компонуется со
слиянием/доверием/устареванием без изменений и --- в отличие от
координатной конвенции, которую оно заменяет, --- знает, когда оно не
знает.

\paragraph{Категорный индикатор: дефект кругового обхода сопряжения.}
Восстановление системы координат --- это в точности та операция,
которая нужна категорному прочтению из \S\ref{sec:categorical}, как
только общий базис $X$ отброшен: при приватных системах координат
каждый агент несёт собственную категорию мест, а коммуникация --- это
пара функторов восстановления $F=\textsc{align}(A,B)$ (система
координат отправителя $B$, выраженная в координатах получателя $A$) и
$G=\textsc{align}(B,A)$. Если они действительно сопряжены --- системы
координат склеены, --- круговой обход $\eta_X: X\to G(F(X))$ есть
тождество, так что \emph{дефект сопряжения}
\[
\mathrm{def}=\lVert F.\mathrm{offset}+G.\mathrm{offset}\rVert_1
\ \text{(сдвиг)},\qquad
r_F{+}r_G\equiv 0\,(\mathrm{mod}\,4)\ \text{и}\
\lVert R_{r_G}F.\delta+G.\delta\rVert_1\ (SE(2)),
\]
равен нулю \citep{maclane1998categories}. Измеренный на описанном выше
методе восстановления (\texttt{experiments/warp/alignment\_defect.py},
3 сдвига $\times$ 10 сидов на ячейку): дефект \emph{в точности} равен
$0$ в $30/30$ запусков для сдвиговых систем координат и в $30/30$ для
каждого из четырёх поворотов на $90^\circ$ --- восстановленные функторы
являются строгими сопряжениями. Там, где восстановление обязано
отказать (агент, наблюдавший только безликую внутреннюю полосу), $F$
или $G$ не определён в $30/30$ запусков с \emph{нулём} ложных склеек
($\mathrm{def}>0$): сопряжения не существует, и выравниватель
отказывается его строить. Точный дискретный выравниватель, таким
образом, реализует дихотомию --- строгое сопряжение
($\mathrm{def}{=}0$) или отсутствие сопряжения --- и никогда не даёт
уверенной неверной склейки; <<отказ закрытым образом>> выше --- это в
точности \emph{отказ от несопряжённой пары}. Это и есть та величина,
которую категорная модель вычисляет, а простое переописание --- нет:
скаляр, равный нулю тогда и только тогда, когда локальные системы
координат коллектива когерентно склеены, и именно эта склейка --- а не
общая координатная конвенция --- была молчаливой предпосылкой слияния
из \S\ref{sec:categorical}.

\section{Категорное BDI-прочтение коллективной семантической памяти}
\label{sec:categorical}

Этот раздел даёт компактное теоретическое прочтение примера поиска
воды в терминах убеждения, желания и намерения (belief, desire,
intention). Цель --- не ввести новый планировщик, а сделать явной
математическую структуру, уже используемую архитектурой коллективной
семантической памяти.

Мы предполагаем конечное множество агентов $A=\{1,\ldots,N\}$, общий
пространственный базис $X$ и общий семантический алфавит $\Sigma$. В
текущей грид-реализации $X$ --- это общая система координат, а
$\Sigma$ содержит теги вроде \texttt{water\_source},
\texttt{water\_candidate}, \texttt{hazard\_edge} и \texttt{open}. Это
явное модельное предположение. Агенты могут иметь разные наблюдения,
разные памяти и разные значения уверенности, но они интерпретируют
места в общей координатной системе и выражают семантические
свидетельства в общем алфавите тегов. Без первого предположения
пиринговое слияние требует сначала решить задачу выравнивания мест ---
раздел~\ref{sec:semantic-identity} решает её конструктивно (консенсус
отпечатков восстанавливает межагентную систему координат), а
раздел~\ref{sec:colimit} формулирует, чем эта конструкция является
категорно. Выравнивание символов (общий $\Sigma$) остаётся
предположением.

У этого предположения есть, кроме того, измеримая \emph{внутренность}:
внутри фиксированного $\Sigma$ вопрос о том, какие классы тегов
фактически несут соответствие мест, --- эмпирическое свойство, а не
модельный выбор. Зарегистрированная абляция словаря в сопутствующей
статье о переносе маршрутов делает сам алфавит воздействием и
обнаруживает, что механизм идентичности \emph{безопасен в $\Sigma$}:
обеднение словаря вплоть до одного содержательного класса превращает
восстановление системы отсчёта в отказ и никогда --- в неверное
выравнивание (ноль fail-open переносов на $252$ ячейках), а
консенсусная масса разложима по классам (опасности $0.56$, стены
$0.44$, границы мира $0.00$). Чего пока не делает ни один механизм ---
это \emph{отбора} $\Sigma$: здесь признаки получают статус ориентира
от дизайна плюс готовности мира сотрудничать. Правило, по которому
признаки \emph{зарабатывают} этот статус участием в успешных
выравниваниях --- алфавит, формируемый координацией, а не
предполагаемый, --- самый острый открытый конец этого предположения и
естественная точка встречи с программой выравнивания символов
(\S\ref{sec:symbol-alignment}).

В момент $t$ агент $i$ поддерживает локальный семантический контекст
$\mathcal{K}_t^i=(G_t^i,\Sigma,I_t^i)$, где $G_t^i$ --- множество
записей о местах в приватной памяти агента, а
\[
I_t^i \subseteq G_t^i \times \Sigma
\]
--- пороговое отношение инцидентности, утверждающее, что место $g\in
G_t^i$ несёт тег $\sigma\in\Sigma$ с достаточной уверенностью.
Градуированная реализация использует уверенности тегов и взвешенные
оценки; чёткое отношение $I_t^i$ --- это пороговый скелет, нужный для
теоретического описания.

Этот формальный контекст индуцирует соответствие Галуа между
множествами тегов и множествами мест:
\[
\operatorname{ext}_t^i(B)=\{g\in G_t^i:\forall\sigma\in B,\ (g,\sigma)\in I_t^i\},
\qquad
\operatorname{int}_t^i(S)=\{\sigma\in\Sigma:\forall g\in S,\ (g,\sigma)\in I_t^i\}.
\]
Для задачи поиска воды желаемый семантический предикат ---
$B_{\mathrm{water}}=\{\texttt{water\_source}\}$ или, в реализованной
взвешенной форме, предикат, комбинирующий положительные свидетельства
для \texttt{water\_source}, \texttt{water\_candidate} и
\texttt{water\_nearby} со штрафами за \texttt{hazard\_edge} или
близкий риск. Извлечение успешно для агента $i$ в момент $t$ ровно
тогда, когда
\[
\operatorname{ext}_t^i(B_{\mathrm{water}})\neq\varnothing .
\]
Таким образом, стадия извлечения --- это не поиск по координатам. Это
проверка того, что запрошенный семантический смысл уже имеет непустой
объём в реализованной памяти агента.

Теперь можно определить компоненты BDI.

\paragraph{Убеждение.} Состояние убеждений агента $i$ --- это
структурированная семантическая память, индуцированная
$\mathcal{K}_t^i$. Эквивалентно, это решётка концептов
$L_t^i=\operatorname{Concepts}(\mathcal{K}_t^i)$, узлы которой
сопрягают объёмы мест с содержаниями тегов. Убеждение о воде --- это
не пропозиция <<вода находится в координате $x$>>. Это область
$L_t^i$, чьё содержание включает водный предикат. Если эта область
пуста, агент не просто плохо выбирает; соответствующий семантический
объект отсутствует в его текущей памяти.

\paragraph{Желание.} Желание --- это обусловленный состоянием
семантический предикат над решёткой убеждений. В примере с водой
внутреннее состояние агента содержит переменную жажды. Активное
желание можно записать как
\[
D_t^i =
\begin{cases}
B_{\mathrm{water}}, & \text{если } \mathrm{thirst}_t^i\geq\theta_{\mathrm{on}},\\
B_{\mathrm{goal}}, & \text{иначе.}
\end{cases}
\]
Итак, желание --- не дополнительный метафизический объект. Это
правило, выбирающее, какой семантический предикат следует запрашивать
из текущей структуры убеждений.

\paragraph{Намерение.} Намерение --- это обязательство по одному
материализованному кандидату. Как только желание выбрало предикат, а
извлечение дало непустой объём, планировщик применяет
удовлетворяющую (satisficing) функцию выбора
\[
\chi_t^i:\mathcal{P}(G_t^i)\rightarrow G_t^i\cup\{\bot\},
\qquad
g_t^{i,\star}=\chi_t^i\!\big(\operatorname{ext}_t^i(D_t^i)\big),
\]
где $\bot$ означает неудачу материализации цели. В терминах Q/R/M/C
желание определяет запрос $Q^\star$, убеждение поддерживает извлечение
$R^\star$, а намерение соответствует материализации $M^\star$.
Завершение $C^\star$ находится ниже по потоку от этой семантической
структуры: оно зависит от того, достигает ли физический контроллер
зафиксированной цели и остаётся ли цель доступной при многоагентной
занятости.

Это даёт компактное категорное прочтение. Пусть $\mathbf{Ctx}_\Sigma$
--- категория, объекты которой суть семантические контексты
$(G,\Sigma,I)$ над фиксированным алфавитом $\Sigma$, а морфизмы
сохраняют семантическую инцидентность. Траектория локальной памяти
агента $i$ --- это индексированная временем диаграмма
\[
\mathcal{K}_0^i\rightarrow\mathcal{K}_1^i\rightarrow\cdots
\rightarrow\mathcal{K}_t^i\rightarrow\mathcal{K}_{t+1}^i\rightarrow\cdots
\]
в $\mathbf{Ctx}_\Sigma$. Каждая стрелка индуцирована наблюдением и
обновлением памяти. Когда новое свидетельство консолидируется,
формальный контекст меняется; следовательно, меняется его решётка
концептов. Таково порядково-теоретическое содержание роста памяти.

В терминах Q/R/M/C пустой объём
$\operatorname{ext}_t^i(D_t^i)=\varnothing$ --- это отказ
убеждения/извлечения, непустой объём без устойчивой фиксации --- отказ
материализации, а успешная фиксация, которую невозможно достичь, ---
отказ завершения: семантическая и физическая части состояния агента
отказывают по-разному и измеримы по отдельности.

Коллективный случай поднимает это на одну конструкцию выше. В моменты
рассылки получатель $i$ держит диаграмму локальных объектов памяти
$\mathcal{D}_t^i=(\mathcal{K}_t^i,\mathcal{K}_t^{j_1},\ldots,
\mathcal{K}_t^{j_m})$ вместе с картами выравнивания и запрашивает
слитый консенсусный объект
$\widetilde{\mathcal{K}}_t^i=\operatorname{Merge}_{\mathrm{trust},
\mathrm{stale}}(\mathcal{D}_t^i)$, где свидетельства соседей
взвешиваются в точности по правилам раздела~\ref{sec:csm}
($\mathrm{trust}_{j\to i}\cdot e^{-\alpha\,\mathrm{age}_j}$). Те же
BDI-определения применимы к слитому контексту: намерение теперь
выбирается из $\operatorname{ext}_{\widetilde{\mathcal{K}}_t^i}(D_t^i)$,
а не из одного лишь приватного объёма.

Это объясняет, почему каденция коммуникации порождает сдвиг узкого
места. Более быстрая рассылка повышает свежесть свидетельств соседей.
Это, как правило, расширяет или улучшает водный объём
$\operatorname{ext}_{\widetilde{\mathcal{K}}_t^i}(B_{\mathrm{water}})$,
так что агенты материализуют цели надёжнее; в терминах Q/R/M/C это
повышает $P(M^\star\mid R^\star)$. Но то же свежее общее убеждение
синхронизирует и намерения. Несколько агентов могут зафиксироваться на
одном и том же дефицитном источнике воды. Память семантически точна,
но физическая цель теперь оспариваема. Это понижает $P(C^\star\mid
M^\star)$. Узкое место, следовательно, физически обосновано. Это не
формальный артефакт декомпозиции: оно возникает потому, что
семантическое убеждение разделяется через память, тогда как завершение
ограничено занятостью и дефицитными ресурсами. Коллективная
семантическая память должна поэтому регулировать не только объём
разделяемой информации, но и то, как она сливается, чьим
свидетельствам доверяют и когда устаревшие свидетельства должны
перестать влиять на формирование намерений.

В этом смысле <<коллектив>> --- не центральный разум и не глобальная
карта, а стабилизированное семейство индексированных соседями
процессов; следующий подраздел делает это утверждение точным.

Двойственное, процессное прочтение --- агенты как коалгебры (codata),
коллектив как финальная коалгебра, коммуникация как корекурсия глубины
один, а сдвиг узкого места как соревнование между бисимилярностью
убеждений и бисимилярностью намерений --- развито в
приложении~\ref{sec:coalgebra}. Это линза, а не механизм: ничто ниже
по тексту от него не зависит, поэтому оно и живёт в приложении.

\subsection{Слияние как копредел и его инварианты}
\label{sec:colimit}

Объект слияния выше можно описать точнее. Работая в версии
$\mathbf{Ctx}_\Sigma$, обогащённой (enrichment) метками времени и
провенансом, слияние снимков от агентов $1,\dots,N$ есть диаграмма
локальных объектов памяти плюс карты выравнивания, индуцированные
отношением тождества мест. Консенсусный объект
$\widetilde{\mathcal{K}}_t^i$ --- это копредел (colimit) этой
диаграммы, конкретно реализуемый слиянием union-find над выровненными
концептами. Такое прочтение даёт три полезных ограничения.
\begin{enumerate}[leftmargin=20pt,itemsep=2pt,topsep=2pt]
\item \textbf{Независимость от порядка.} Слитый объект единствен с точностью до изоморфизма, поэтому результат не является артефактом обработки снимков агентов в конкретном порядке.
\item \textbf{Цепочки ожидаемы.} Если концепт A выравнивается с B, а B --- с C, union-find выполняет транзитивное замыкание. Чрезмерное слияние, следовательно, не баг структуры данных; это цена выбранного отношения эквивалентности.
\item \textbf{Затухание --- не морфизм.} Устаревание хранится как поле веса вне морфизмов выравнивания. Иначе <<состарить-потом-слить>> и <<слить-потом-состарить>> не обязаны коммутировать, что разрушило бы чистую интерпретацию широковещательных снимков как одной диаграммы.
\end{enumerate}

В ранних версиях этот раздел заканчивался признанием, что модель не
отвечает на вопрос, какие перцептивные признаки должны отождествлять
два места --- карты выравнивания предполагались, а не строились. Этот
пробел теперь закрыт разделом~\ref{sec:semantic-identity}, и это
закрытие даёт копредельному прочтению его недостающую научную функцию:
\begin{enumerate}[leftmargin=20pt,itemsep=2pt,topsep=2pt]
\item \textbf{Морфизмы выравнивания строятся, а не предполагаются.}
  Консенсус отпечатков И ЕСТЬ алгоритм, порождающий карты выравнивания
  диаграммы --- взаимно уникальные сопоставления ориентиров определяют
  преобразование системы координат, через которое вкладывается каждый
  локальный объект памяти.
\item \textbf{Существование становится эмпирическим предикатом.}
  Копредел существует тогда и только тогда, когда существует
  консенсус. В $180^\circ$-симметричном мире морфизм выравнивания не
  существует \emph{в принципе} (две несравнимые диаграммы
  поддерживают равный консенсус), и реализация обнаруживает именно
  это и отказывается сливать --- поведение <<отказа закрытым образом>>
  из раздела~\ref{sec:semantic-identity} есть операционализированное
  категорное несуществование. Этот предикат измеряется, а не
  постулируется: дефект кругового обхода сопряжения из
  \S\ref{sec:semantic-identity} в точности равен $0$, когда попарное
  выравнивание --- строгое сопряжение (во всех восстановленных
  случаях сдвига и поворота на $90^\circ$, $30/30$), и не определён,
  когда сопряжения не существует (отказ закрытым образом, $30/30$, с
  нулём ложных склеек) --- скалярный свидетель существования
  копредела.
\item \textbf{Затухание остаётся вне морфизмов} (инвариант 3), и
  эмпирически вентиль доверие$\times$устаревание не меняется при
  восстановлении системы координат --- те же целочисленные точки
  излома с общей системой координат и без неё, --- что подтверждает
  коммутирование двух слоёв, как того требует диаграммное прочтение.
\end{enumerate}

\subsection{Коллективная память как когерентная склейка локальных категорий мест}
\label{sec:gluing}

Копредельное слияние из \S\ref{sec:colimit} \emph{предполагало} карты
выравнивания, а \S\ref{sec:semantic-identity} их \emph{строит}; теперь мы можем
сформулировать в качестве центрального категорного результата, каким объектом
является коллективная память, когда агенты держат \emph{разные} системы
координат, и в точности когда она существует точным (faithful) образом.

\paragraph{Постановка.} Реализованная память агента $i$ --- это тонкая
категория (частично упорядоченное множество): её решётка концептов
$L_i=\operatorname{Concepts}(\mathcal{K}_i)$ над приватной системой координат
$i$, объекты --- места, порядок --- включение объёмов из
\S\ref{sec:categorical}. Восстановление системы координат
(\S\ref{sec:semantic-identity}) даёт на сонаблюдаемом подпосете монотонное
отображение $F_{ij}:L_i\to L_j$, переносящее места $i$ в систему координат
$j$, и обратное к нему $F_{ji}$. Поскольку категории \emph{тонкие},
восстановленная пара $(F_{ij},F_{ji})$ является \emph{сопряжённой парой} ровно
тогда, когда она есть соответствие Галуа, с единицей
$\eta^{ij}:\operatorname{Id}_{L_i}\Rightarrow F_{ji}\!\circ\!F_{ij}$ ---
дефектом сопряжения, измеренным в \S\ref{sec:semantic-identity}.

\paragraph{Определение (когерентная система выравнивания).} Семейство
$\{F_{ij}\}$ над графом коммуникации \emph{когерентно}, если (i) каждая пара
--- сопряжённое порядковое вложение ($\eta^{ij}$ --- изоморфизм, т.е.\ дефект
$0$: ни одно место не вклеивается в другое), и (ii) оно удовлетворяет условию
коцикла $F_{jk}\!\circ\!F_{ij}=F_{ik}$ на тройных перекрытиях.

\paragraph{Предложение (два агента).} Для сопряжённой пары $F\dashv G$ между
$L_A,L_B$ категория-коллаж (запятая) $\mathrm{Gl}(F)$ --- объекты
$L_A\sqcup L_B$, перекрёстные морфизмы $\mathrm{Gl}(F)(X,Y)=L_B(FX,Y)\cong L_A(X,GY)$
\citep{maclane1998categories} --- содержит $L_A,L_B$ как полные подпосеты и
является двухагентной коллективной памятью; дефект $\eta_X$ --- препятствие к
тому, чтобы $X$ вернулся в себя круговым обходом, и склейка отождествляет
различные места тогда и только тогда, когда $\eta$ не изоморфизм.

\paragraph{Теорема (коллективная память = копредел диаграммы выравнивания).}
Пусть $\mathbf{D}$ --- диаграмма в $\mathbf{Pos}$ (эквивалентно, в обогащённой
$\mathbf{Ctx}_\Sigma$), вершины которой --- $L_i$, а рёбра --- $F_{ij}$.
Коллективная память есть $\operatorname{colim}\mathbf{D}$. Так как
$\mathbf{Pos}$ кополна, этот копредел всегда существует; он является
\emph{точной} склейкой --- каждая $L_i$ вкладывается и никакие два различных
места не отождествляются --- \emph{тогда и только тогда}, когда система
выравнивания когерентна. Когда все $F_{ij}=\operatorname{id}$ (общая система
координат), $\mathbf{D}$ --- постоянная диаграмма и
$\operatorname{colim}\mathbf{D}$ --- в точности копредельное union-find-слияние
из \S\ref{sec:colimit}; общее утверждение устраняет предположение об общей
системе координат, заменяя тождества восстановленными сопряжёнными вложениями.

\emph{Аргумент.} Кополнота $\mathbf{Pos}$ реализует копредел как дизъюнктное
объединение, факторизованное по предпорядку, порождённому
$F_{ij}$-отождествлениями. Каждая каноническая инъекция
$L_i\hookrightarrow\operatorname{colim}\mathbf{D}$ точна тогда и только тогда,
когда эти отождествления не добавляют схлопываний внутри $L_i$, что имеет
место тогда и только тогда, когда каждое $F_{ij}$ --- порядковое вложение
(единица --- изоморфизм, дефект $0$) и отождествления согласованы на
перекрытиях (коцикл) --- т.е.\ когерентность. Случай общей системы координат
--- постоянная диаграмма, чей копредел есть покомпонентное объединение, т.е.\
union-find. \hfill$\square$

\paragraph{Следствие (два вычислимых свидетеля).} Когерентность проверяема без
решения общей задачи семантического выравнивания. (a)~Попарный \emph{дефект
сопряжения} $\eta$ (\S\ref{sec:semantic-identity}) равен нулю тогда и только
тогда, когда каждое ребро --- строгое сопряжённое вложение --- измерен в
точности $0$ в $30/30$ восстановленных случаев сдвига и поворота --- и
защищает от \emph{чрезмерного слияния}. (b)~\emph{Связность} графа
достижимости доверие$\times$каденция (\S\ref{sec:phase-diagram},
рис.~\ref{fig:phase-diagram}) решает, покрывает ли один
$\operatorname{colim}\mathbf{D}$ всех $N$ или диаграмма распадается на
копроизведение под-консенсусов (\emph{фрагментация}). Эти два механизма отказа
категорно различны и измеряются по отдельности: дефект $>0$ --- это
схлопывание, несвязность --- фрагментация: две оси
рис.~\ref{fig:phase-diagram} и поведение отказа закрытым образом из
\S\ref{sec:semantic-identity}.

\paragraph{Границы применимости.} Доказательства --- стандартная кополнота
$\mathbf{Pos}$ и конструкция коллажа \citep{maclane1998categories}; работа
теоремы в том, чтобы назвать объект и привязать его существование к двум
измеренным величинам. Эмпирическое содержание --- что восстановление по
отпечаткам действительно даёт сопряжённые вложения --- живёт в
\S\ref{sec:semantic-identity}, а не в теореме. Одно предположение \emph{не}
снято: общий алфавит тегов $\Sigma$; при разных алфавитах $F_{ij}$ должны были
бы переводить ещё и $\Sigma$, что выходит за пределы текущего восстановления и
является честной следующей границей.

\subsection{Параметры протокола как категорная структура: фазовая диаграмма доверие\texorpdfstring{$\times$}{x}каденция}
\label{sec:phase-diagram}

Прочтение до сих пор моделировало \emph{память} (объекты --- семантические
контексты, слияние --- их копредел) при общей системе координат и полном
обмене. Его предсказательная польза приходит, когда мы направляем линзу на
\emph{структуру коммуникации} и спрашиваем, \emph{какую форму принимает
копредел} как функцию двух параметров протокола --- чтобы категория вычисляла
структурный факт, а не переописывала алгоритм. (В $\mathbf{Pos}$ копредел
диаграммы \emph{существует} всегда --- по кополноте, связна диаграмма или
нет; структурный вопрос в том, является ли он единой консенсусной
компонентой или копроизведением под-консенсусов --- в точности дихотомия
\S\ref{sec:gluing}.) Два параметра играют две
структурные роли:
\begin{description}[leftmargin=12pt,topsep=2pt,itemsep=1pt]
\item[Доверие $\tau$ управляет морфизмами.] Ребро $i\!\to\!j$ (агент $i$
включает свидетельства $j$) существует только там, где доверие $i$ к $j$ не
ниже $\tau$. Повышение $\tau$ прореживает граф доверия.
\item[Каденция $K$ управляет глубиной композиции.] На горизонте $H$ есть
$R=\lfloor H/K\rfloor$ раундов рассылки, поэтому составная цепочка
$i\!\to\!k\!\to\!\cdots\!\to\!j$ может иметь длину не более $R$ морфизмов.
Низкое $K$ (много раундов) допускает глубокую композицию; высокое $K$ ---
только короткие пути.
\end{description}
\emph{Охватывающий} (consensus-spanning) копредел --- единая связная
консенсусная компонента, сплавляющая и обслуживающая свидетельства каждого
агента --- получается тогда и только тогда, когда $R$-раундовая достижимость
$\tau$-прореженного графа доверия является все-ко-всем; иначе копредел
вырождается в копроизведение под-консенсусов.

Рисунок~\ref{fig:phase-diagram} перебирает $\tau\times K$ для $N{=}8$ агентов,
чьё попарное доверие порождено собственным EMA-правилом системы
(\method{peer\_memory.peer\_trust.AsymmetricTrust}) при кольцевой топологии
сонаблюдения (соседи сонаблюдают часто и зарабатывают высокое доверие; далёкие
агенты --- нет), усреднённой по $12$ сидам
(\method{experiments/collective\_semantic\_memory/phase\_diagram\_trust\_cadence.py}).
Результат --- настоящий фазовый переход с диагональной границей. При низком
$\tau$ граф доверия плотен, и копредел охватывает всех агентов при любой
каденции. Повышение $\tau$ прореживает дальние морфизмы до почти-соседского
кольца; коллектив тогда сохраняет охватывающий копредел \emph{только если}
каденция достаточно низка, чтобы многошаговая композиция через всё ещё
доверенных посредников вновь его связала --- $\tau{=}0.55$ удерживает
охватывающий копредел вплоть до $R{=}2$,
$\tau{=}0.80$ требует $R\ge4$ (диаметр кольца), а $\tau{=}0.90$ фрагментирован
повсюду. Каденция и доверие, следовательно, не непрозрачные ручки: вместе они
определяют категорный структурный инвариант, и фазовая граница --- это в
точности место, где коллектив переходит от единого консенсусного объекта к
фрагментированным под-консенсусам.

\begin{figure}[t]
  \centering
  \includegraphics[width=\linewidth]{figures/fig_phase_diagram.pdf}
  \caption{Фазовая диаграмма структуры коммуникации по порогу доверия $\tau$
  (строки) и каденции $K$ (столбцы; $R=\lfloor H/K\rfloor$ раундов рассылки),
  $N{=}8$ агентов на кольце взаимодействия, $12$ сидов. \textbf{Слева:}
  плотность достижимости (доля упорядоченных пар агентов, чьи свидетельства
  могут скомпоноваться за $R$ раундов). \textbf{Справа:}
  $P(\text{копредел охватывает всех }N)$, т.е.\ достижимость
  все-ко-всем. Диагональная граница --- структурный переход: коллектив
  допускает единый консенсусный объект в области низких $\tau$/низких $K$ и
  фрагментируется в углу высоких $\tau$/высоких $K$, причём более высокое
  $\tau$ требует более низкого $K$ (больше раундов), чтобы остаться связным.}
  \label{fig:phase-diagram}
\end{figure}

Именно в этом смысле категорная модель операциональна, а не описательна:
каждый параметр протокола отображается в конкретное изменение диаграммы
коммуникации, а структурный инвариант диаграммы (охватывает ли копредел всех
агентов или распадается в копроизведение) предсказывает качественный режим
коллектива. \emph{Границы применимости.}
Матрица доверия порождена настоящим EMA-правилом при смоделированной
кольцевой топологии сонаблюдения, а <<охватывающий консенсус>>
операционализирован как $R$-раундовая достижимость все-ко-всем; количественно
установить, что пересечение этой структурной границы предсказывает падение
наблюдаемой Q/R/M/C-производительности --- естественный следующий эксперимент.

\subsection{К выравниванию символов: заземление алфавитов тегов через общие места}
\label{sec:symbol-alignment}

Теорема о склейке (\S\ref{sec:gluing}) всё ещё предполагает общий алфавит
$\Sigma$. Отбрасывая его, каждый агент держит контекст $(G_i,\Sigma_i,I_i)$
над \emph{собственным} словарём, и правильный морфизм --- уже не функтор над
одним $\Sigma$, а \emph{инфоморфизм} \citep{barwise1997information}: пара
$f:G_A\!\to\!G_B$, $\hat f:\Sigma_B\!\to\!\Sigma_A$ с условием
$I_B(f g,\,y)\Leftrightarrow I_A(g,\,\hat f y)$. Инфоморфизмы компонуются и
имеют копределы (каналы Барвайза--Селигмана), поэтому коллективная память над
гетерогенными алфавитами --- это \emph{ядро канала} --- копредел диаграммы
инфоморфизмов, --- что дословно расширяет \S\ref{sec:gluing} (функторы над
одним $\Sigma$ становятся инфоморфизмами над $\Sigma_i$; общий $\Sigma$ с
$\hat f=\mathrm{id}$ --- прежний случай).

\paragraph{Механизм заземления.} Символы выравниваются через общие референты:
на местах, уже заякоренных восстановлением системы координат
(\S\ref{sec:semantic-identity}), совместная встречаемость $A$-тега и $B$-тега
--- свидетельство того, что они именуют один и тот же признак, и $\hat f$
считывается с заякоренной совместной встречаемости. (Когда алфавиты
различаются настолько, что сами отпечатки не могут посеять якоря, начальная
инициализация идёт от алфавитно-независимых структурных констелляций с
чередованием восстановления $f$ и $\hat f$.) \emph{Дефект выравнивания
символов} $\mathrm{def}_\Sigma$ --- доля заякоренных пар
$(\text{место},y)$, нарушающих бикондиционал инфоморфизма, --- символьный
аналог дефекта сопряжения систем координат из \S\ref{sec:semantic-identity}.

\paragraph{Предварительный результат.} Контролируемое исследование
(\method{experiments/warp/symbol\_alignment.py}; $|\Sigma_A|{=}8$, $120$ мест,
$20$ сидов) даёт агенту $B$ секретную переразметку алфавита $A$ с шумом
наблюдения и восстанавливает $\hat f$ из заякоренной совместной
встречаемости.

\begin{table}[h]
\caption{Выравнивание символов заземлением против предположения об идентичных
алфавитах. Target-ID F1 --- идентификация целевого признака $A$ по
(переведённым) свидетельствам $B$. <<naive>> предполагает совпадение алфавитов
по индексу; <<recovered>> выучивает $\hat f$; <<oracle>> --- истинное
отображение.}
\label{tab:symbol-alignment}
\centering
\small
\begin{tabular}{c c c c c c}
\toprule
шум & точн.\ перевода & $\mathrm{def}_\Sigma$ & F1 naive & F1 recovered & F1 oracle \\
\midrule
0.00 & 1.00 & 0.000 & 0.56 & 1.00 & 1.00 \\
0.10 & 1.00 & 0.099 & 0.55 & 0.92 & 0.92 \\
0.30 & 0.89 & 0.304 & 0.56 & 0.69 & 0.73 \\
\bottomrule
\end{tabular}
\end{table}

Предположение об идентичных алфавитах проваливается
($\text{F1}\approx0.55$: переразметка отправляет целевой тег $B$ на неверный
признак); заземление восстанавливает отображение (точность перевода $1.00$,
$\mathrm{def}_\Sigma$ отслеживает уровень шума) и достигает уровня оракула в
идентификации цели, устойчиво вплоть до $15\%$ заякоренных мест (точность
$0.95$). Огрубление двух признаков в один $B$-тег ограничивает достижимую
точность --- точной биекции не существует, поэтому ядро канала огрубляет, а
не ошибается. Выравнивание символов, таким образом, не полностью открытая
проблема, а тот же рецепт <<заземление плюс дефект>> на один алфавит выше;
полное рассмотрение (подлинная несоизмеримость, онлайн-дрейф) --- предмет
следующей статьи.

\section{Ограничения и следующие эксперименты}
\label{sec:limitations}

\textbf{Грид-субстрат.} Текущий результат намеренно контролируем: все
эксперименты проведены в грид-мирах. Субстрат с непрерывным состоянием, такой
как VMAS, Habitat, ProcTHOR или MiniWorld, --- следующий шаг переносимости.

\textbf{Координатный якорь.} Устранён в
разделе~\ref{sec:semantic-identity}: локальные сигнатуры-отпечатки
точно восстанавливают межагентную систему координат (сдвиг и повороты
на $90^\circ$), и результаты по слиянию/доверию/устареванию не
меняются поверх восстановленных систем. Открытым остаётся более
трудный конец той же оси: непрерывное $SE(2)$, шум внешнего вида
(текущий сопоставитель использует точное равенство констелляций,
поскольку субстрат бесшумен). Выравнивание символов --- агенты с
\emph{разными алфавитами тегов} --- предварительно рассмотрено в
\S\ref{sec:symbol-alignment} (заземление через общие места,
восстановление в виде инфоморфизма); подлинная несоизмеримость и дрейф
словаря остаются открытыми.

\textbf{Минимальная CSM.} Правила CSM намеренно просты. Адаптивная
каденция $K_t$ и доверие, обусловленное намерением, --- следующие
архитектурные варианты. У селективной консолидации недопредставленных
тегов теперь есть измеренный пилот: в сопутствующей статье о переносе
маршрутов чисто реляционная \emph{песня} (сигнатуры ориентиров $+$
такт-векторы, без глобальных координат) выполняет тот же перенос
никогда не виденной цели, что и полный снапшот, за $6.4\%$ битовой
стоимости, причём безопасно --- одни такты отказывают открытым
образом, одни сигнатуры не достигают невиданного. Замена обмена
снапшотами на бюджетированные песни, сливаемые теми же тремя правилами
CSM, перестала быть спекуляцией; открытой остаётся вторая половина ---
сделать грамматику песни (что триггерит куплет, что входит в
сигнатуру) предметом отбора, а не дизайна.

\textbf{Обучаемые базовые варианты.} Обучаемая коммуникация в духе
CommNet --- полезное сравнение, но более релевантный будущий базовый
вариант --- многоагентная система с явной LLM-памятью или векторной
памятью, измеренная на той же задаче.

\section{Заключение}

Коллективная память не монотонна. В дефицитной многоагентной
навигации мгновенный обмен информацией может снизить завершение,
заставляя агентов сходиться на одних и тех же целях. Пиринговая
каденция обнажает этот компромисс, но одна лишь каденция --- одномерная
поверхность управления. Коллективная семантическая память должна
сделать слияние, доверие и устаревание явными. Минимальная CSM делает
ровно это и переводит систему в новую рабочую точку: не насыщение M
при быстрой рассылке и не насыщение C при медленной, а более
сбалансированный профиль M$\times$C с более высоким успехом, чем у
каждой фиксированной пиринговой каденции на бенчмарке дефицита.

\begin{ack}
Рукопись подготовлена в анонимизированной форме для рецензирования.
Описанные эксперименты выполнены на том же бенчмарк-конвейере, что и в
сопутствующей статье об измерениях Q/R/M/C, но архитектурное
утверждение здесь ограничено коллективной семантической памятью.
\end{ack}

{\small
\begin{thebibliography}{9}

\bibitem[Aczel(1988)]{aczel1988nonwellfounded}
P.~Aczel.
\newblock \emph{Non-Well-Founded Sets}.
\newblock CSLI Lecture Notes 14, 1988.

\bibitem[Barwise and Seligman(1997)]{barwise1997information}
J.~Barwise and J.~Seligman.
\newblock \emph{Information Flow: The Logic of Distributed Systems}.
\newblock Cambridge University Press, 1997.

\bibitem[Bettini et~al.(2022)]{bettini2022vmas}
M.~Bettini, R.~Kortvelesy, J.~Blumenkamp, and A.~Prorok.
\newblock VMAS: A vectorized multi-agent simulator for collective robot learning.
\newblock In \emph{16th International Symposium on Distributed Autonomous Robotic Systems (DARS)}, 2022.
\newblock arXiv:2207.03530.

\bibitem[Boyd et~al.(2006)]{boyd2006gossip}
S.~Boyd, A.~Ghosh, B.~Prabhakar, and D.~Shah.
\newblock Randomized gossip algorithms.
\newblock \emph{IEEE Trans.\ Information Theory}, 52(6):2508--2530, 2006.

\bibitem[Chevalier-Boisvert et~al.(2023)]{MinigridMiniworld23}
M.~Chevalier-Boisvert, B.~Lahlou, S.~Willems, and P.~S. Castro.
\newblock MiniGrid and MiniWorld: Modular reinforcement learning environments.
\newblock \emph{arXiv:2306.13831}, 2023.

\bibitem[Deitke et~al.(2022)]{deitke2022procthor}
M.~Deitke, E.~V. Wijmans, D.~Schwenk, O.~Salvador, K.~M. Ehsani, et al.
\newblock ProcTHOR: Large-scale embodied AI using procedural generation.
\newblock In \emph{NeurIPS Datasets and Benchmarks}, 2022.

\bibitem[Foerster et~al.(2018)]{foerster2018ctde}
J.~N. Foerster, G.~Farquhar, T.~Afouras, N.~Nardelli, and S.~Whiteson.
\newblock Counterfactual multi-agent policy gradients.
\newblock In \emph{AAAI}, 2018.

\bibitem[Jacobs(2016)]{jacobs2016coalgebra}
B.~Jacobs.
\newblock \emph{Introduction to Coalgebra: Towards Mathematics of States and Observation}.
\newblock Cambridge University Press, 2016.

\bibitem[Mac~Lane(1998)]{maclane1998categories}
S.~Mac~Lane.
\newblock \emph{Categories for the Working Mathematician}.
\newblock Springer, 2nd edition, 1998.

\bibitem[Rao and Georgeff(1995)]{rao1995bdi}
A.~S. Rao and M.~P. Georgeff.
\newblock BDI agents: From theory to practice.
\newblock In \emph{ICMAS}, 1995.

\bibitem[Rutten(2000)]{rutten2000universal}
J.~J.~M.~M. Rutten.
\newblock Universal coalgebra: a theory of systems.
\newblock \emph{Theoretical Computer Science}, 249(1):3--80, 2000.

\bibitem[Sangiorgi(2011)]{sangiorgi2011bisimulation}
D.~Sangiorgi.
\newblock \emph{Introduction to Bisimulation and Coinduction}.
\newblock Cambridge University Press, 2011.

\bibitem[Savva et~al.(2019)]{savva2019habitat}
M.~Savva, A.~Kadian, O.~Maksymets, Y.~Zhao, E.~Wijmans, B.~Jain, J.~Straub, J.~Liu, V.~Koltun, J.~Malik, et al.
\newblock Habitat: A platform for embodied AI research.
\newblock In \emph{ICCV}, 2019.

\bibitem[Sperber and Wilson(1995)]{sperber1995relevance}
D.~Sperber and D.~Wilson.
\newblock \emph{Relevance: Communication and Cognition}.
\newblock Blackwell, 2nd edition, 1995.

\bibitem[Sukhbaatar et~al.(2016)]{sukhbaatar2016commnet}
S.~Sukhbaatar, A.~Szlam, and R.~Fergus.
\newblock Learning multiagent communication with backpropagation.
\newblock In \emph{NeurIPS}, 2016.

\bibitem[Xiao et~al.(2007)]{xiao2007distributed}
L.~Xiao, S.~Boyd, and S.~Lall.
\newblock A scheme for robust distributed sensor fusion based on average consensus.
\newblock \emph{IPSN}, 2007.

\end{thebibliography}
}

\appendix

\section{Коалгебраический слой: индивиды как ко-данные, коллектив как финальная коалгебра}
\label{sec:coalgebra}

Прочтение из \S\ref{sec:categorical} фиксирует \emph{содержание}
убеждений: в данном состоянии решётка концептов
$L_t^i=\operatorname{Concepts}(\mathcal{K}_t^i)$ говорит, что извлекаемо,
а слияние из \S\ref{sec:colimit} --- как содержания соседей комбинируются
в $\widetilde{\mathcal K}_t^i$. Оба \emph{алгебраичны}: они строят объект
убеждения. Ни одно из них не говорит об агенте как о незавершающемся
\emph{процессе}, который обязан выдавать действие и сигнал на каждом
такте в движущейся среде, и ни одно --- о том смысле, в котором
\emph{популяция} приходит к общей конвенции. Эта процессная сторона ---
\emph{коалгебраическая} (coalgebra)~\citep{rutten2000universal,jacobs2016coalgebra}.
Мы делаем её точной на примере поиска воды и показываем, что две стороны
двойственны и встречаются внутри одного перехода состояния. (Словарь
убеждение/желание/намерение следует традиции BDI-агентов
\citep{rao1995bdi}, здесь наделённой коалгебраическим содержанием.)

\paragraph{Сигнатура.} Зафиксируем множество действий $\mathrm{Act}$
(двигаться, копать, ничего не делать), множество наблюдений $O_i$ на
агента (локальный перцепт агента $i$ --- его \emph{уникальная точка
зрения}) и множество \emph{остензий}
\[
  M \;=\; 1 \,+\, (X\times\Sigma),
\]
значения которого --- либо \emph{молчание} (слагаемое $1$), либо
указатель <<место $x$ несёт тег $\sigma$>> (свистнуть и указать на
клетку с меткой воды). То, что молчание --- полноправное значение $M$,
существенно: остензивный канал \emph{может быть информативно пуст}, и он
отделён от обновления состояния ниже. Это ровно та наблюдаемость,
которую разрушает обучаемый вещественный вектор сообщения (его канал
никогда не пуст и никогда не отделяет сигнализацию от шага), и именно
поэтому символьные стадии остаются читаемыми.

\paragraph{Состояние.} В качестве носителя коалгебры возьмём
высокоуровневое состояние H-MDP из \S\ref{sec:categorical},
\[
  \mathcal{S}_i \;=\; \big\{\, h=(\mathcal{K}^i,\,D^i,\,g^{i,\star},\,s^i) \,\big\},
\]
несущее ту же информацию, что и $(L^i,D^i,g^{i,\star},s^i)$, поскольку
$L^i=\operatorname{Concepts}(\mathcal K^i)$. Таким образом,
\emph{убеждение живёт в состоянии}; решётка --- его проекция.

\paragraph{Определение (ко-данные: котип агента).}
Наблюдаемое поведение агента $i$ --- это коалгебра для эндофунктора
\[
  B_i(S)\;=\;\underbrace{\mathrm{Act}}_{\text{действовать сейчас}}\;\times\;
            \underbrace{M}_{\text{сигналить сейчас}}\;\times\;
            \underbrace{S^{\,O_i\times M^{A\setminus\{i\}}}}_{\substack{\text{продолжить при (мой следующий}\\ \text{перцепт, сигналы соседей)}}},
\]
т.е.\ отображение $c_i=\big\langle\, \method{observe\_act}_i,\ \method{observe\_ostension}_i,\ \method{step\_intension}_i \,\big\rangle : \mathcal{S}_i \to B_i(\mathcal{S}_i)$ с тремя деструкторами
\[
  \method{observe\_act}_i:\mathcal{S}_i\to\mathrm{Act},\quad
  \method{observe\_ostension}_i:\mathcal{S}_i\to M,\quad
  \method{step\_intension}_i:\mathcal{S}_i\to \mathcal{S}_i^{\,O_i\times M^{A\setminus\{i\}}} .
\]
Здесь $M^{A\setminus\{i\}}=\prod_{j\neq i}M$ --- кортеж остензий соседей,
доставленных на этом такте. Операционально
$\method{step\_intension}_i(h)(o,(m_j)_{j\neq i})$
(i)~вкладывает перцепт $o$ и указатели соседей $(m_j)$ в контекст через
слияние с доверием/устареванием,
$\widetilde{\mathcal K}_t^i=\operatorname{Merge}_{\mathrm{trust,stale}}(\cdots)$
из \S\ref{sec:colimit}; (ii)~пересчитывает решётку; (iii)~переоценивает
предикат желания $D^i$; и (iv)~перефиксирует намерение
$g^{i,\star}=\chi_t^i(\operatorname{ext}(D^i))$ из \S\ref{sec:categorical}.
Итак, $\operatorname{Merge}$ и $\chi$ --- \emph{компоненты одного
деструктора}. Агент \emph{и есть} ко-данные, порождаемые $c_i$: у него
нет конечного терминального плана, лишь постоянная способность выдавать
действие и остензию и продолжать. Это <<бесконечный генератор
выживания>> из мотивирующего примера.

\paragraph{Определение (индексированное семейство).}
Система --- это $A$-индексированное семейство
$\{(\mathcal{S}_i,c_i)\}_{i\in A}$, коалгебра в $\mathbf{Set}^{A}$.
Функтор $B_i$ зависит от индекса $i$ через $O_i$ (и через множество
соседей $A\setminus\{i\}$): один и тот же котип выживания
инстанцируется у каждого агента через локальную точку зрения этого
агента. Агент~1 на дюне порождает из <<я вижу горизонт>>; агент~2 в
низине порождает из <<я чувствую почву>>. Индекс --- это точка зрения.

\paragraph{Определение (козависимость: коллективная коалгебра).}
Пусть среда доставляет совместный перцепт $o=(o_i)_{i\in A}\in\prod_i O_i$.
Коллективное однотактовое отображение на $\prod_i\mathcal{S}_i$
подключает выход остензии каждого агента ко входам шага остальных,
\[
  \Phi\big((h_i)_i\big)(o)\;=\;
  \Big(\ \method{step\_intension}_i(h_i)\big(o_i,\,(\method{observe\_ostension}_j(h_j))_{j\neq i}\big)\ \Big)_{i\in A}.
\]
Поскольку следующее состояние $i$ есть функция \emph{выходов
деструкторов} остальных, семейство --- не $N$ независимых процессов, а
одна коалгебра для функтора-произведения
$\mathbf{B}(S)=\big(\prod_i(\mathrm{Act}\times M)\big)\times S^{\prod_i O_i}$
на $\prod_i\mathcal{S}_i$. Эта проводка и есть \emph{козависимость}:
агент~1 не может сделать шаг, не пропустив остензию агента~2 через свой
собственный $\method{step\_intension}$. Каденция рассылки входит сюда
как вентиль: на тактах без рассылки доставляемый кортеж $(m_j)$ --- это
последний полученный (устаревший) кортеж или молчание; $K$ --- период
доставки.

\paragraph{Определение (поведение и бисимуляция).}
$B_i$ --- контейнерный (полиномиальный) функтор, поэтому у него есть
финальная коалгебра
$(\nu B_i,\ \mathrm{beh}_i)$~\citep{aczel1988nonwellfounded,rutten2000universal},
элементы которой --- поведенческие деревья с ветвлением по
$\big(O_i\times M^{A\setminus\{i\}}\big)$ и метками
$(\mathrm{Act}\times M)$, и существует единственный морфизм коалгебр
$!_i:\mathcal{S}_i\to\nu B_i$, отправляющий состояние во всё его будущее
поведение. \emph{$B_i$-бисимуляция} (bisimulation) --- это отношение
$R\subseteq\mathcal{S}_i\times\mathcal{S}_i$, при котором связанные
состояния выдают одинаковое действие и одинаковую остензию и на каждом
входе переходят в $R$-связанные состояния. Контейнерные функторы
сохраняют слабые обратные образы, поэтому \emph{наибольшая бисимуляция}
$\sim_i$ существует, является эквивалентностью и равна $\ker(!_i)$: два
состояния бисимилярны тогда и только тогда, когда у них одинаковое
поведение~\citep{sangiorgi2011bisimulation}. То же верно для
коллективной коалгебры
$(\prod_i\mathcal{S}_i,\ \langle(\method{observe\_act}_i)_i,(\method{observe\_ostension}_i)_i,\Phi\rangle)$
с наибольшей бисимуляцией $\approx$.

\paragraph{Определение (общество / культура).}
\emph{Коллектив} --- это образ состояний в финальной коалгебре
$\nu\mathbf{B}$, эквивалентно фактор
$\big(\prod_i\mathcal{S}_i\big)/\!\approx$. Это ни центральный сервер,
ни общая карта: это класс бисимилярности совместного поведения.
\emph{Конвенция поиска воды} --- этот инвариант, прочитанный на водной
подсигнатуре. Пусть $M_{\mathrm{water}}\subseteq M$ --- указатели, чей
тег лежит в водном семействе, и пусть
$\nu B_i\to\nu B^{\mathrm{w}}$ --- забывающее отображение в
поведенческие деревья, ограниченные приёмом $M_{\mathrm{water}}$ и
выдачей соответствующей фиксации намерения. Популяция \emph{держит
конвенцию}, когда достижимые поведения всех агентов согласуются под
этим отображением: услышав <<укажи-и-копай у отмеченного водой места>>,
деструктор намерения каждого агента фиксирует это место (с точностью до
занятости). Управление тогда автоматично в обоих направлениях. Новый
агент (четвёртый агент из примера) --- член общества ровно тогда, когда
его поведение попадает в класс $[\cdot]_{\approx}$ на
$M_{\mathrm{water}}$; иначе его остензии не понимаются и его собственный
шаг никогда не сходится (сверху-вниз); а сам класс был достигнут
локальными, связанными шагами первых агентов у одного зелёного куста
(снизу-вверх).

\paragraph{Замечание (коммуникация = ленивая корекурсия глубины один).}
Поведение $!_i(h)$ определяется бесконечным развёртыванием $c_i$.
Получив $m_j=\method{observe\_ostension}_j(h_j)$, агент $i$ использует
внутри $\method{step\_intension}_i$ только этот выход деструктора
\emph{глубины один}; он \emph{не} вычисляет $!_j(h_j)\in\nu B_j$ ---
полное поведение соседа, каковое и есть регресс <<я думаю, что ты
думаешь\ldots>>. Коммуникация --- это, следовательно, корекурсия,
усечённая до наименьшей глубины, разрешающей собственный следующий шаг
получателя, --- конечное приближение отображения в финальную коалгебру
--- с релевантностно-теоретическим правилом остановки
\citep{sperber1995relevance}, решающим эту глубину. Регресс схлопывается
в одно эффективное действие.

\paragraph{Замечание (мост алгебра--коалгебра к \S\ref{sec:categorical}--\ref{sec:colimit}).}
Убеждение алгебраично: FCA-решётка $L_t^i$, построенная копредельным
слиянием, --- это \emph{содержание в состоянии}. Агент коалгебраичен:
ко-данные, наблюдаемые тремя деструкторами, --- это \emph{процесс над
состояниями}. Они двойственны и встречаются в одном месте ---
копредельное слияние и удовлетворяющий выбор $\chi$ суть внутренности
$\method{step\_intension}$. Диаграмма роста памяти
$\mathcal{K}_0^i\to\mathcal{K}_1^i\to\cdots$ из \S\ref{sec:categorical}
--- это в точности проекция-на-убеждения траектории коалгебры
$h_0^i\mapsto h_1^i\mapsto\cdots$.

\paragraph{Коалгебраическое прочтение сдвига узкого места.}
Результат \S\ref{sec:cadence} и таблицы~\ref{tab:csm-results} читается
здесь одной строкой. Соревнуются два вида бисимилярности.
Бисимилярность на подсигнатуре \emph{остензий / убеждений} --- желанный
инвариант: это общая водная конвенция. Бисимилярность на сыром выходе
\emph{намерений} над \emph{дефицитной} целью --- нежеланный: несколько
состояний, отображающихся в одно и то же $g^{\star}$, --- это коллизия.
Занятость $\mathcal{O}_{-i}$ --- внекоалгебраическое ограничение,
штрафующее второе. Каденция $K$ настраивает, насколько сильно связка
$\Phi$ гонит коллектив к \emph{полной} наибольшей бисимуляции
$\approx$: малое $K$ доставляет свежие остензии часто, расширяя слитый
водный объём, так что намерение фиксируется надёжно
($P(M^\star\mid R^\star)\uparrow$), но пересинхронизирует деструктор
намерения, так что агенты сталкиваются
($P(C^\star\mid M^\star)\downarrow$). Внутренний оптимум $K{=}8$ ---
это точка, где убеждения сходятся, а намерения не схлопываются на одну
клетку; минимальная CSM добавляет доверие и устаревание именно затем,
чтобы $\method{step\_intension}$ достигал бисимилярности убеждений
(свежие, доверенные, выровненные свидетельства), гася бисимилярность
намерений (устаревшие свидетельства затухают прежде, чем успеют согнать
всех агентов на одну и ту же фиксацию). Мы формулируем это как
прочтение, согласующееся с данными, а не как теорему: контейнерные
функторы, их финальные коалгебры и бисимуляция --- стандартные
конструкции, используемые здесь как линза, а
релевантностно-теоретическое правило остановки --- интерпретирующая
глосса, а не подогнанный механизм.

\begin{table}[t]
\caption{Термины мотивирующего примера как формальные объекты
коалгебраического прочтения. Каждая строка --- смысл конструкции, уже
использованной в \S\ref{sec:coalgebra}, в задаче поиска воды.}
\label{tab:coalg-dictionary}
\centering
\small
\begin{tabular}{p{2.9cm} p{4.3cm} p{4.4cm}}
\toprule
Неформальный термин & Формальный объект & Смысл в поиске воды \\
\midrule
Котип (ко-данные) & $B_i$-коалгебра $(\mathcal{S}_i,c_i)$ & незавершающийся генератор действий/сигналов агента \\
Индексированный тип & семейство над $A$; зависимость $B_i$ от $O_i$ & уникальная локальная точка зрения агента (дюна против низины) \\
Козависимый тип & проводка $\method{observe\_ostension}_j\!\to\!\method{step\_intension}_i$ в $\Phi$ & <<я шагаю только через твою остензию>> \\
Индивид (экземпляр) & состояние $h\in\mathcal{S}_i$ и его поведение $!_i(h)$ & работающий агент, копающий и свистящий сейчас \\
Деструкторы (когенераторы) & $\method{observe\_act}$, $\method{observe\_ostension}$, $\method{step\_intension}$ & шагни-и-копай; свистни-и-укажи; вложи сигнал соседа в следующее состояние \\
Ленивое вычисление & корекурсия глубины один + релевантностная остановка & развернуть коалгебру соседа на один шаг, а не всё его поведение \\
Финальное коиндуктивное ядро / общество & $\nu\mathbf{B}$; наибольшая бисимуляция $\approx$ (на $M_{\mathrm{water}}$) & стабильная распределённая культура поиска воды \\
\bottomrule
\end{tabular}
\end{table}


\section{Команды воспроизведения}
\label{app:repro}

\begin{itemize}[leftmargin=16pt]
\item Основная развёртка по каденции:
\path{python experiments/big_experiment/exp_cadence_phase.py --mode full --workers 8 --out_dir tmp/big_experiment_qrmc_40}
\item Расширенные каденции:
\path{python experiments/big_experiment/exp_extra_K.py --workers 8 --out_dir tmp/big_experiment_extraK}
\item Анализ Q/R/M/C:
\path{python experiments/big_experiment/analyze_qrmc.py --runs_csv tmp/big_experiment_qrmc_40/runs.csv --out_dir tmp/big_experiment_qrmc_40}
\item Бенчмарк минимальной CSM:
\path{python -m experiments.collective_semantic_memory.run_csm_benchmark --out_dir tmp/csm_benchmark}
\item Развёртка переносимости на MiniGrid:
\path{python -m experiments.minigrid_multiagent_wrapper.run_portability_sweep --out_dir tmp/minigrid_portability}
\item Дефект кругового обхода сопряжения при восстановлении системы координат (раздел~\ref{sec:semantic-identity}):
\path{python experiments/warp/alignment_defect.py}
\item Структурная фазовая диаграмма доверие$\times$каденция (раздел~\ref{sec:phase-diagram}, рис.~\ref{fig:phase-diagram}):
\path{python experiments/collective_semantic_memory/phase_diagram_trust_cadence.py --fig figures/fig_phase_diagram.pdf}
\end{itemize}

\section{Замечание об активах и лицензиях}
\label{app:assets}

В экспериментах используются MiniGrid и MiniWorld~\citep{MinigridMiniworld23}. Локально реализованный код бенчмарка и памяти будет выпущен под лицензией MIT в составе артефакт-пакета.

\clearpage
\FloatBarrier
\input{checklist.tex}

\end{document}
